\documentclass[UTF8,zihao=-4]{ctexart}
\usepackage[a4paper,margin=2.2cm]{geometry}
\usepackage{xcolor}
\usepackage{hyperref}
\usepackage{enumitem}
\usepackage{fontspec}
\usepackage{amsmath}
\IfFontExistsTF{Microsoft YaHei}{\setCJKmainfont{Microsoft YaHei}}{\setCJKmainfont{SimSun}[BoldFont=SimHei]}
\IfFontExistsTF{Microsoft YaHei UI}{\setCJKsansfont{Microsoft YaHei UI}}{}
\IfFontExistsTF{Consolas}{\setmonofont{Consolas}}{}
\definecolor{linkblue}{RGB}{30,80,145}
\hypersetup{colorlinks=true,linkcolor=linkblue,urlcolor=linkblue,pdfborder={0 0 0}}
\setlength{\parindent}{2em}
\setlength{\parskip}{0.25em}
\linespread{1.16}
\setlist{itemsep=0.2em,topsep=0.35em,leftmargin=2.5em}
\pagestyle{plain}

\title{08\_线性时间选择算法}

\author{}

\date{}

\begin{document}

\maketitle

\section*{08 线性时间选择算法}

\subsection*{题目原文}

线性时间选择算法的实现和应用

\subsection*{考查类型}

算法设计题和复杂度分析题。

\subsection*{问题定义}

选择问题是：给定 $n$ 个元素和整数 \texttt{k}，找出第 \texttt{k} 小的元素。

特殊情况：

\begin{quote}
\small\ttfamily\raggedright
\relax k\ =\ 1\ \ \ \ \ \ \ \ 找最小值\\
\relax k\ =\ n\ \ \ \ \ \ \ \ 找最大值\\
\relax k\ =\ (n+1)/2\ \ 找中位数\\
\end{quote}

\subsection*{随机选择算法思想}

随机选择与快速排序的划分类似：

\begin{enumerate}
\item 选一个划分基准。
\item 按基准把数组分成小于、等于、大于三部分。
\item 根据 \texttt{k} 所在范围，只递归进入一边。
\end{enumerate}

平均时间复杂度：

\begin{quote}
\small\ttfamily\raggedright
\relax O(n)\\
\end{quote}

最坏时间复杂度：

\[O(n^{2})\]

\subsection*{确定性线性时间选择算法}

也称中位数的中位数方法。

步骤：

\begin{enumerate}
\item 将 $n$ 个元素每 5 个分成一组。
\item 对每组内部排序，取每组中位数。
\item 递归地从这些中位数中选出中位数，作为划分基准。
\item 用该基准对原数组划分。
\item 根据 \texttt{k} 的位置，只递归进入一边。
\end{enumerate}

\subsection*{为什么是线性时间}

每组 5 个元素。取所有组的中位数后，再取这些中位数的中位数作为划分基准。

至少一半的组中位数不小于该基准。每个这样的组中，至少有 3 个元素不小于该组中位数，因此也不小于该基准。取整项只影响常数，不影响线性时间结论。

由此可保证每次递归至少排除常数比例的元素。

递推可写为：

\[T(n) \le  T(n/5) + T(7n/10) + O(n)\]

解得：

\[T(n) = O(n)\]

\subsection*{应用}

找中位数：使用 $Select(A, \lceil n/2 \rceil)$。

找第 $k$ 小：

\begin{quote}
\small\ttfamily\raggedright
\relax Select(A,\ k)\\
\end{quote}

找前 \texttt{k} 小元素：

先用选择算法找到第 \texttt{k} 小的元素，再按该元素划分数组。

找分位点：

递归使用选择算法找多个秩位置的元素。

\subsection*{伪代码框架}

\begin{quote}
\small\ttfamily\raggedright
\relax Select(A,\ k):\\
\relax \ \ \ \ if\ |A|\ is\ small:\\
\relax \ \ \ \ \ \ \ \ sort\ A\\
\relax \ \ \ \ \ \ \ \ return\ A[k]\\
\mbox{}\\
\relax \ \ \ \ split\ A\ into\ groups\ of\ 5\\
\end{quote}
\[M = medians of all groups\]
\[p = Select(M, \lceil |M| / 2 \rceil)\]
\begin{quote}
\small\ttfamily\raggedright
\mbox{}\\
\relax \ \ \ \ partition\ A\ into\ L,\ E,\ G\ by\ p\\
\mbox{}\\
\relax \ \ \ \ if\ k\ ≤\ |L|:\\
\relax \ \ \ \ \ \ \ \ return\ Select(L,\ k)\\
\relax \ \ \ \ else\ if\ k\ ≤\ |L|\ +\ |E|:\\
\relax \ \ \ \ \ \ \ \ return\ p\\
\relax \ \ \ \ else:\\
\relax \ \ \ \ \ \ \ \ return\ Select(G,\ k\ -\ |L|\ -\ |E|)\\
\end{quote}

\subsection*{例题}

\subsubsection*{例题 1}

题目：在数组中找第 \texttt{5} 小元素：

\begin{quote}
\small\ttfamily\raggedright
\relax A\ =\ [9,\ 1,\ 7,\ 3,\ 8,\ 2,\ 6,\ 5,\ 4]\\
\end{quote}

解答：

按升序排列用于核对结果：

\begin{quote}
\small\ttfamily\raggedright
\relax [1,\ 2,\ 3,\ 4,\ 5,\ 6,\ 7,\ 8,\ 9]\\
\end{quote}

第 \texttt{5} 小元素为：

\begin{quote}
\small\ttfamily\raggedright
\relax 5\\
\end{quote}

选择算法不需要完整排序。一次以 \texttt{5} 为基准划分：

\begin{quote}
\small\ttfamily\raggedright
\relax L\ =\ [1,\ 3,\ 2,\ 4]\\
\relax E\ =\ [5]\\
\relax G\ =\ [9,\ 7,\ 8,\ 6]\\
\end{quote}

因为：

\[|L| = 4\]
\[|L| + |E| = 5\]

第 \texttt{5} 小元素正好落在 $E$ 中，答案是 \texttt{5}。

\subsubsection*{例题 2}

题目：对下面数组使用中位数的中位数方法选取划分基准，并求第 \texttt{8} 小元素。

\begin{quote}
\small\ttfamily\raggedright
\relax [12,\ 5,\ 7,\ 1,\ 9,\ 3,\ 14,\ 6,\ 10,\ 2,\ 8,\ 15,\ 4,\ 13,\ 11]\\
\end{quote}

解答：

每 5 个一组：

\begin{quote}
\small\ttfamily\raggedright
\relax [12,\ 5,\ 7,\ 1,\ 9]\ \ \ \ \ \ \ 中位数\ 7\\
\relax [3,\ 14,\ 6,\ 10,\ 2]\ \ \ \ \ \ 中位数\ 6\\
\relax [8,\ 15,\ 4,\ 13,\ 11]\ \ \ \ \ 中位数\ 11\\
\end{quote}

中位数组为：

\begin{quote}
\small\ttfamily\raggedright
\relax [7,\ 6,\ 11]\\
\end{quote}

其中位数是 \texttt{7}，取 \texttt{7} 为划分基准。

划分后：

\begin{quote}
\small\ttfamily\raggedright
\relax L\ =\ [5,\ 1,\ 3,\ 6,\ 2,\ 4]\\
\relax E\ =\ [7]\\
\relax G\ =\ [12,\ 9,\ 14,\ 10,\ 8,\ 15,\ 13,\ 11]\\
\end{quote}

因为：

\[|L| = 6\]
\[|L| + |E| = 7\]

第 \texttt{8} 小元素在 \texttt{G} 中，并且是 \texttt{G} 的第 \texttt{1} 小元素。\texttt{G} 的最小值为：

\begin{quote}
\small\ttfamily\raggedright
\relax 8\\
\end{quote}

因此原数组第 \texttt{8} 小元素是 \texttt{8}。

\subsection*{常见误区}

线性时间选择算法不需要完整排序。完整排序是 $O(n \log n)$，选择算法只递归处理一边。

每组大小取 5 是为了保证递推能收敛到线性时间。组大小不是随意选择的。

\end{document}
